%----------------------------------------------------------------------------------------
%	CONCLUSION
%----------------------------------------------------------------------------------------

\section{Conclusion}

This work presented a comprehensive framework for the control of a double pendulum with nonlinear joint stiffness, addressing the \textbf{challenges of under-actuation and chaotic dynamics}.

Regarding trajectory generation, we showed that standard \textbf{Newton-based approaches} are limited in \textbf{complex scenarios}. The introduction of the \textbf{Dynamic Horizon Extension} proved crucial for generating \textbf{intricate feasible paths} that not only reach the target but maintain terminal stability.

Finally, we implemented and compared two tracking controllers: a Linear Time-Varying Linear Quadratic Regulator (LTV-LQR), and a Model Predictive Control (MPC). The comparison revealed that \textbf{MPC} is essential when operating \textbf{near actuator limits}, as it explicitly handles constraints that the LTV-LQR cannot.

Conversely, the \textbf{LTV-LQR} demonstrated superior performance in terms of \textbf{computational efficiency and ease of implementation} for less demanding tasks.

Overall, the results were visualized using a custom simulation tool, confirming that the proposed framework can successfully generate and track complex trajectories.

\newpage